% ===========================================================================
% OWNER: M1 (Zoya Ahmad)
% GRADED BY: Form 3 #7
% DRAW FROM: evidence/README.md, evidence/extract.py, docs/04-Implementation/evidence-collector.md, docs/evidence-m1.md, ADR-002, ADR-006
% Rules: mid-eval-report/CLAUDE.md. Do not edit a file you do not own.
% ===========================================================================
\subsection{Resource Evidence Extraction (Member 1)}\label{ch6:sec:m1}
Member 1's module is the \path{evidence/} package. Given a polluter and a victim found by
Member 2's runner, it names the shared resource the polluter writes and the victim reads, and the
exact bytecode locations of that write and read. It works only from the project's compiled
classes, read with the JDK's \texttt{javap}: it runs no tests, adds no instrumentation and never
writes to the project (ADR-001, ADR-002). It is 716 lines of Python in one file, using only the
standard library, with 505 lines of tests (\texttt{wc -l}, 12 October 2026). Figure~\ref{ch5:fig:evidence-dfd}
shows its data flow.

\subsubsection{Files and main functions}
{\small
\begin{longtable}{>{\raggedright\arraybackslash}p{0.27\textwidth}>{\raggedright\arraybackslash}p{0.07\textwidth}>{\raggedright\arraybackslash}p{0.54\textwidth}}
\caption{Files of the resource evidence module}\label{ch6:tab:m1-files}\\
\toprule
File & Lines & Job\\
\midrule
\endfirsthead
\toprule
File & Lines & Job\\
\midrule
\endhead
\path{extract.py} & 716 & The extractor and its command line; main functions below\\
\path{tests/test_extract.py} & 505 & 39 \texttt{unittest} tests in seven classes, run on real
  \texttt{javap} output\\
\path{tests/resources/m1-selftest/} & -- & A small Maven project of ten Java files written as test
  input (lifecycle methods, call depth, JUnit~3, unsupported cases); not part of the fixture\\
\path{tools/measure_depth.py} & 70 & Measures accesses, unsupported observations, \texttt{javap}
  calls and time per depth; produced the depth data behind ADR-006\\
\path{javap-dumps/} & -- & Saved \texttt{javap} output quoted as evidence in
  \path{docs/evidence-m1.md}\\
\bottomrule
\end{longtable}}

The main functions in \path{extract.py}, in the order a pair is analysed:
\begin{itemize}
\item \texttt{Project.load}: runs \texttt{javap -v -p -c} on one class and parses it into
  \texttt{ClassInfo}, \texttt{Method} and \texttt{Instruction} objects, cached per class.
\item \texttt{find\_roots}: the test method plus its lifecycle methods (\texttt{@Before},
  \texttt{@After}, \texttt{@BeforeClass}, \texttt{@AfterClass}, JUnit~3 \texttt{setUp} and
  \texttt{tearDown}) and static initialisers, collected from the test class and its project
  superclasses; a method overridden in a subclass is skipped, because it never runs.
\item \texttt{walk\_root}, \texttt{scan\_method}, \texttt{scan\_call}, \texttt{follow\_call}: a
  breadth-first walk from each root that classifies each instruction and follows calls into
  project classes up to the depth limit.
\item \texttt{analyse\_test}: all accesses, unsupported observations and library calls not
  followed, for one test.
\item \texttt{find\_edges}: the resources the polluter writes and the victim reads.
\item \texttt{analyse\_pair}: both tests and their edges, deepening when needed.
\item \texttt{report\_fields}: turns a pair into the three report fields and the limitations.
\item \texttt{main}: the command line, \texttt{python3 -m evidence.extract}, which prints JSON and
  exits 0 (done), 1 (\texttt{javap} missing or failed) or 2 (input error).
\end{itemize}

\subsubsection{How extraction works}
For each root, \texttt{walk\_root} scans the method's bytecode one instruction at a time and queues
the project methods it calls, breadth-first, so every access is reported at the shortest call depth
at which it can be reached. Each method is scanned at most once per root, which also stops recursion
and call cycles. Listing~\ref{ch6:lst:m1-scan} shows the core of the scan.

\begin{lstlisting}[language=Python,basicstyle=\footnotesize\ttfamily,columns=fullflexible,keepspaces=true,caption={Core of \texttt{scan\_method}, abridged (\texttt{evidence/extract.py})},label=ch6:lst:m1-scan]
for index, ins in enumerate(method.instructions):
    owner = ins.owner or info.name
    if ins.opcode in ("getstatic", "putstatic"):
        declaring = project.declaring_class(owner, ins.name)
        resource = {"kind": "static-field",
                    "class": declaring, "field": ins.name}
        access = "READ" if ins.opcode == "getstatic" else "WRITE"
        result["accesses"].append(make_access(
            "static-field", resource, access, walk["via"],
            info, method, ins.offset, walk["path"]))
    elif ins.opcode == "invokedynamic":
        result["unsupported"].append(observation(
            "INVOKEDYNAMIC", info, method, ins, ...))
    elif ins.opcode.startswith("invoke"):
        scan_call(project, chain_names, info, method,
                  index, owner, walk, result)
\end{lstlisting}

Each instruction is classified as follows:
\begin{itemize}
\item \texttt{getstatic} and \texttt{putstatic} are a READ and a WRITE of a static field. The field is
  named by the class that \emph{declares} it, so a field reached through a subclass is still one
  resource.
\item A call to \texttt{System.getProperty}, \texttt{Integer.getInteger}, \texttt{Long.getLong} or
  \texttt{Boolean.getBoolean} is a READ of a system property. A call to
  \texttt{System.setProperty} or \texttt{clearProperty} is a WRITE, but only when every argument is pushed by one simple instruction
  right before the call and the first is a constant string (\texttt{constant\_key}); otherwise
  the call is reported as unsupported.
\item A call into a project class is followed one level deeper (\texttt{follow\_call}). A call into
  the JDK or a library is listed in \texttt{external\_calls\_not\_followed}.
\item Anything the extractor cannot model becomes a typed \emph{unsupported observation} with its
  location, never a guess: \texttt{REFLECTION}, \texttt{INVOKEDYNAMIC},
  \texttt{VIRTUAL\_\allowbreak DISPATCH} (only the named target is followed, not subclass overrides),
  \texttt{UNRESOLVED\_\allowbreak CALL}, \texttt{SYSPROP\_\allowbreak NON\_\allowbreak CONSTANT\_\allowbreak KEY},
  \texttt{SYSPROP\_\allowbreak BULK}, \texttt{IMPLICIT\_\allowbreak CLINIT} (a static initialiser
  that may run on first use), \texttt{DEPTH\_\allowbreak LIMIT} and
  \texttt{JUNIT\_\allowbreak RULE\_\allowbreak OR\_\allowbreak RUNNER}.
\end{itemize}
Every access records the resource, READ or WRITE, the class, method and bytecode offset, the root
it was reached from (\texttt{via}, for example \texttt{TEST\_METHOD} or \texttt{BEFORE}) and the call
path. \texttt{find\_edges} then keeps each resource the polluter writes and the victim reads, with
every write and read location. Writes made by the victim itself are not edges.

\subsubsection{Call depth and auto-deepening}\label{ch6:sec:m1-depth}
The depth is how many calls deep the walk follows from a root: depth~1 scans only the root's own
bytecode, depth~2 also the project methods it calls, and so on. Each extra level finds more but costs
more time and adds more code that is reachable but may never run. The team first fixed the default at
2 with a maximum of 3 (ADR-002). Measurements on the two real fastjson cases from the proof of concept
showed that this was not enough (Table~\ref{ch6:tab:m1-depth}): the victims reach the shared field
only inside constructors, four or five calls deep.

{\small
\begin{longtable}{>{\raggedright\arraybackslash}p{0.10\textwidth}>{\raggedright\arraybackslash}p{0.12\textwidth}>{\raggedright\arraybackslash}p{0.22\textwidth}>{\raggedright\arraybackslash}p{0.22\textwidth}>{\raggedright\arraybackslash}p{0.18\textwidth}}
\caption{Edges found per depth on fastjson (\texttt{docs/evidence-m1.md}, 10 October 2026)}\label{ch6:tab:m1-depth}\\
\toprule
Case & Depth & Real candidates with an edge & Other candidates with an edge & Seconds, all candidates\\
\midrule
\endfirsthead
\toprule
Case & Depth & Real candidates with an edge & Other candidates with an edge & Seconds, all candidates\\
\midrule
\endhead
FJ-01 & 1--3 & 0 of 42 & 0 of 30 & 22.4--30.2\\
FJ-01 & 4 & 42 of 42 & 1 of 30 & 34.8\\
FJ-01 & 5 & 42 of 42 & 1 of 30 & 45.8\\
FJ-02 & 1--3 & 0 of 24 & 0 of 698 & 143.3--156.1\\
FJ-02 & 4 & 1 of 24 & 0 of 698 & 163.2\\
FJ-02 & 5 & 24 of 24 & 0 of 698 & 180.7\\
\bottomrule
\end{longtable}}

A ``real'' candidate is one that changed the victim's result when the proof of concept ran it first.
Raising the fixed depth to~5 would make every pair pay the deepest cost: one pair took 3--6~s at
depth~2, but 7.5~s for FJ-01 at depth~4 and 26.3~s for FJ-02 at depth~5 (ADR-006). The team therefore
adopted \emph{auto-deepening} (ADR-006), implemented in \texttt{analyse\_pair}:
\begin{enumerate}
\item analyse both tests at the requested depth (default~2);
\item if there is no edge \emph{and} a walk stopped at the depth limit (\texttt{DEPTH\_LIMIT}),
  repeat one level deeper;
\item stop at the first depth with an edge, when nothing was cut off, or at depth~5.
\end{enumerate}
The pair records \texttt{depth\_requested} and \texttt{depth\_used}, and when the reported evidence
is deeper than the default, \texttt{report\_fields} adds a limitation saying that deeper evidence is
more likely to be code that never runs. On the fixture, every ground-truth edge is found at depth~2
without deepening.

\subsubsection{How the runner uses the module}
Member 2's \texttt{diagnose} command calls the module once the polluters are known
(\texttt{resource\_fields} in \path{runner/cli.py}): it runs \texttt{analyse\_pair} once for each
polluter with the victim, and \texttt{combine\_fields} joins the results. With one polluter the
report holds \texttt{report\_fields} unchanged. With several, it shows the first polluter's edge and
names the others in \texttt{limitations}, but only if \emph{every} polluter has an edge; otherwise the
three resource fields are \texttt{null}, so that a \texttt{VERIFIED} result always means every
blamed test has a found mechanism. The report holds one resource and one location per side, so
\texttt{report\_fields} shows the first edge and its first write and read, and lists any other edges
and locations in \texttt{limitations}. The module can also be run on its own, for example on F2:
\begin{lstlisting}[basicstyle=\footnotesize\ttfamily,columns=fullflexible,keepspaces=true,breaklines=true]
python3 -m evidence.extract --classes fixtures/od-fixture/target/classes \
    --test-classes fixtures/od-fixture/target/test-classes \
    --polluter odfixture.FeaturePolluterTest#enableTurbo \
    --victim odfixture.FeatureVictimTest#expectsTurboDisabled
\end{lstlisting}
Run on 12 October 2026 (main \texttt{82df9fa}), it exited 0 and its \texttt{pair} section reported
one edge, \texttt{sysprop:odfixture.turbo}: written at \texttt{FeaturePolluterTest.enableTurbo}
offset~4 and read at \texttt{FeatureFlags.isTurboEnabled} offset~2, a method the victim calls, so
the read is at depth~2 (\texttt{depth\_used}~2).

\subsubsection{Setup and running}
The module needs Python~3 and a JDK~8 or newer with \texttt{javap} on the path (or
\texttt{FLAKETRACE\_JAVAP} set to another \texttt{javap}); it has no other dependencies. Its tests
need the fixture and the self-test project compiled first. From the repository root:
\begin{lstlisting}[basicstyle=\footnotesize\ttfamily,columns=fullflexible,keepspaces=true,breaklines=true]
mvn -B -q -f fixtures/od-fixture/pom.xml test-compile
mvn -B -q -f evidence/tests/resources/m1-selftest/pom.xml test-compile
FLAKETRACE_REQUIRE_JVM=1 python3 -m unittest -v evidence.tests.test_extract
\end{lstlisting}
\texttt{FLAKETRACE\_REQUIRE\_JVM=1} turns a missing \texttt{javap} into a failure instead of skipped
tests. On Windows Command Prompt the variable is set with \texttt{set FLAKETRACE\_REQUIRE\_JVM=1} and
Python is started with \texttt{py}. These are the same commands as the \texttt{evidence} job in
\path{.github/workflows/ci.yml}, which runs them on Ubuntu with JDK~8 for every pull request.

On 12 October 2026, on main \texttt{82df9fa}, both builds exited 0 and the suite reported
\texttt{Ran 39 tests in 60.392s \ldots{} OK}. This run was made with Claude in the team's cloud
environment (Linux, Maven 3.9.11, OpenJDK and \texttt{javap} 21.0.12.1, Python 3.13.16), not on a
member's own machine; it is recorded in \path{docs/evidence-m1.md}.

\subsubsection{Results}
All results are recorded in \path{docs/evidence-m1.md}.
\begin{itemize}
\item \textbf{Fixture against its ground truth.} Every fixture test is tried as polluter against
  every other as victim: 240 ordered pairs over 16 tests. Edges appear for exactly the ground-truth
  pairs, each on the ground-truth resource, at the default depth without deepening: F1
  (\texttt{Config.mode}), F2 (\texttt{odfixture.turbo}, read at depth~2), F3 (\texttt{Toggles.flagA}
  and \texttt{flagB}, one per polluter) and F4 (\texttt{LateFlag.isSet}). No test gets an edge to N1,
  N2 or N3. The ground truth was written by Member~3 before any run.
\item \textbf{End to end.} The same edges appear in the reports of Table~\ref{ch8:tab:fixture-results},
  where all four order-dependent cases are \texttt{VERIFIED}.
\item \textbf{A real project.} On fastjson at the proof of concept's pinned commit, built unmodified,
  the edges for FJ-01 and FJ-02 appear at depths~4 and~5 (Table~\ref{ch6:tab:m1-depth}), on the
  fields \texttt{JSON.defaultTimeZone} and \texttt{JSON.defaultLocale}; the offsets on FJ-01's read path were checked
  against the saved \texttt{javap} output (\path{evidence/javap-dumps/phase5-fastjson.txt}). Both cases are brittle tests whose earlier test is a state-setter,
  which the module treats the same way: one test writes, the other reads.
\end{itemize}

\subsubsection{Testing}
The 39 tests in \path{evidence/tests/test_extract.py} run the real extractor on real
\texttt{javap} output; nothing is mocked. They are grouped in seven classes:
\texttt{DepthOneExtractionTest} (the exact accesses and offsets, JUnit~3 and~4 lifecycle
attribution, unsupported cases reported, not guessed), \texttt{CallDepthTest} (what each depth adds,
the depth limit, recursion), \texttt{PairSelfTest} (edge rules, including writes in \texttt{setUp} and
static initialisers), \texttt{DeepeningTest}, \texttt{FixturePairTest} (F1--F3 and the report-field
shape), \texttt{GroundTruthTest} (the 240-pair check above) and \texttt{CommandLineTest} (output and
exit codes). When the ground-truth check was added, it was itself checked: forcing depth~1 made it
fail on the missing F2 pair, and altering the ground truth in memory made it fail as well.

\subsubsection{Limitations and a failure case}
Every report carries the module's seven fixed limitations (Section~\ref{ch1:sec:constraints}). Three
matter most in practice:
\begin{itemize}
\item \textbf{A false edge.} On FJ-01 at depths~4 and~5, one of the 30 candidates that did not affect
  the victim still gets an edge: \texttt{DateTest2\#test\_date} writes \texttt{JSON.defaultTimeZone},
  but sets it and restores the old value in \texttt{tearDown}. The extractor sees a WRITE because it
  does not model the values written. This is why an edge alone never makes a result
  \texttt{VERIFIED}: the polluter must also be confirmed by reruns.
\item \textbf{Two edges, one cause.} On both fastjson cases the state-setter's \texttt{setUp} writes
  both \texttt{defaultTimeZone} and \texttt{defaultLocale}, and the victim reads both, so both are
  edges. Static analysis cannot say which one causes the failure; from the test's own constants it
  is the time zone, but that has not been confirmed by a run. The report shows the first edge and
  names the other in \texttt{limitations}.
\item \textbf{Evidence from one real project.} Two real cases from fastjson, a library built around
  global static configuration, are favourable to this kind of analysis. Whether depth~5 is enough
  elsewhere is not yet known; where static analysis finds nothing, the result is
  \texttt{UNRESOLVED} with \texttt{NO\_SUPPORTED\_RESOURCE\_EVIDENCE}, never ``independent'', and
  runtime evidence is the planned fallback (ADR-001).
\end{itemize}
